\chapter{Symbolic State and Process Trees}
\label{chap:symbolic_state}

\section{The Symbolic State Representation}
\label{sec:symbolic_state:struct}

At the heart of the NumLang metacomputation engine lies the \texttt{SymbolicState} struct (\texttt{src/mir/supercompiler/state.rs}). Unlike traditional symbolic execution tools which maintain heavyweight first-order logical constraints for solver-based path queries, NumLang models symbolic execution states as lightweight, canonicalized, SSA-aligned algebraic records:

\begin{lstlisting}[language=Rust, caption={The SymbolicState structure in NumLang.}]
pub struct SymbolicState {
    pub env: HashMap<Place, SymTermId>,
    pub heap: HashMap<SymTermId, SymTermId>,
    pub path_conditions: Vec<PathCondition>,
    pub intervals: HashMap<Place, Interval>,
    pub current_block: BasicBlockId,
    pub call_depth: usize,
    pub loop_depth: usize,
}
\end{lstlisting}

\subsection{Path Conditions and Refinements}
Each \texttt{PathCondition} represents a boolean symbolic fact acquired when driving through conditional branch terminators (\texttt{Terminator::BranchIf} or \texttt{Terminator::Switch}). Path conditions are maintained in conjunctive normal form:
\begin{equation}
\Phi = \bigwedge_{i=1}^m \psi_i, \quad \text{where } \psi_i \in \{\mathtt{BinOp}(\mathtt{Cmp}, t_1, t_2), \mathtt{Discriminant}(v) == k\}
\end{equation}
Crucially, path conditions immediately update the \texttt{intervals} domain (Chapter~\ref{chap:refinements}), enabling instantaneous branch pruning without requiring external SMT invocations during the driving loop.

\section{The SymTerm Intermediate Term Language}
\label{sec:symbolic_state:term_lang}

Symbolic values are modeled by the inductively defined \texttt{SymTerm} language (\texttt{src/mir/supercompiler/term.rs}). A symbolic term represents either a known static scalar, a symbolic input variable, an algebraic combination of terms, an abstract heap reference, or a suspended codata thunk:

\begin{lstlisting}[language=Rust, caption={The SymTerm algebraic datatype.}]
pub enum SymTerm {
    ConstInt(i64, Type),
    ConstFloat(u64, Type),
    ConstBool(bool),
    ConstStr(String),
    Var(Place, Type),
    Binary(BinaryOp, SymTermId, SymTermId, Type),
    Unary(UnaryOp, SymTermId, Type),
    Constructor(String, usize, Vec<SymTermId>, Type),
    Select(SymTermId, SymTermId, SymTermId, Type),
    Phi(Vec<(BasicBlockId, SymTermId)>, Type),
    Call(String, Vec<SymTermId>, Type),
    Ref(SymTermId, Type),
    Deref(SymTermId, Type),
    Discriminant(SymTermId, Type),
    ClosureVal(String, Vec<SymTermId>, Type),
    Thunk(String, Vec<SymTermId>, Type),
}
\end{lstlisting}

\subsection{Small-Step Term Reduction Semantics}
The small-step operational semantics of $\symterm$ evaluation under concrete valuation $\sigma : \mathcal{V} \to \mathcal{C}$ are given by the deterministic reduction relation $t \downarrow v$:
\begin{align}
\text{\bf (E-Const)} \quad & \mathtt{ConstInt}(c, \tau) \downarrow c \\[0.8em]
\text{\bf (E-Var)} \quad & \frac{\sigma(x) = v}{\mathtt{Var}(x, \tau) \downarrow v} \\[0.8em]
\text{\bf (E-BinOp)} \quad & \frac{t_1 \downarrow v_1 \quad t_2 \downarrow v_2 \quad v = v_1 \;\text{op}\; v_2}{\mathtt{Binary}(\text{op}, t_1, t_2, \tau) \downarrow v} \\[0.8em]
\text{\bf (E-Select)} \quad & \frac{t_{\text{cond}} \downarrow \mathtt{true} \quad t_{\text{then}} \downarrow v}{\mathtt{Select}(t_{\text{cond}}, t_{\text{then}}, t_{\text{else}}, \tau) \downarrow v} \qquad
\frac{t_{\text{cond}} \downarrow \mathtt{false} \quad t_{\text{else}} \downarrow v}{\mathtt{Select}(t_{\text{cond}}, t_{\text{then}}, t_{\text{else}}, \tau) \downarrow v}
\end{align}

\section{Hash-Consing via the TermInterner}
\label{sec:symbolic_state:interner}

To ensure that symbolic equivalence checks run in $\mathcal{O}(1)$ time throughout process tree exploration, NumLang implements global \emph{hash-consing} via the \texttt{TermInterner} (\texttt{src/mir/supercompiler/term.rs}). Every unique symbolic term is allocated exactly once and assigned a dense, integer-typed \texttt{SymTermId}.

\subsection{Fast Hashing and Invariants}
Hashing uses a specialized 64-bit rotating hash function (\texttt{fx\_hash\_step}) parameterized by the golden ratio multiplier:
\begin{equation}
h_{k+1} = (h_k \lll 5) \oplus (v \cdot \mathtt{0x517cc1b727220a95})
\end{equation}
The \texttt{TermInterner} maintains the following structural invariants:
\begin{enumerate}
    \item \textbf{Unique Representation}: For any two terms $t_1, t_2$, $\text{lookup}(t_1) = \text{lookup}(t_2) \iff \text{structurally\_identical}(t_1, t_2)$.
    \item \textbf{Monotonic Sizing}: The size of an interned term $t = f(c_1, \dots, c_k)$ is strictly computed as $1 + \sum \text{size}(c_i)$, cached in an array indexable by \texttt{SymTermId}.
    \item \textbf{Subterm Depth}: The depth is cached as $1 + \max \text{depth}(c_i)$, enabling instant whistle size filtering without tree traversal.
\end{enumerate}

\section{The 30 Algebraic Reduction Identities (Phase 59)}
\label{sec:symbolic_state:identities}

During term interning (\texttt{intern\_binary} and \texttt{intern\_unary}), NumLang executes canonical algebraic identity reductions. Rather than generating bloated compound terms that must later be simplified by peephole passes, the interner algebraically normalizes expressions on the fly. 

Table~\ref{tab:algebraic_identities} enumerates the 30 foundational identities implemented in Phase 59.

\begin{table}[h!]
\centering
\small
\begin{tabular}{clll}
\toprule
\textbf{\#} & \textbf{Algebraic Form} & \textbf{Normalized Output} & \textbf{Mathematical Rationale} \\
\midrule
1 & $x + 0$ & $x$ & Additive identity \\
2 & $0 + x$ & $x$ & Commutative additive identity \\
3 & $x - 0$ & $x$ & Subtractive right identity \\
4 & $x - x$ & $0$ & Additive inverse nilpotence \\
5 & $0 - x$ & $-x$ & Negation introduction \\
6 & $x \times 1$ & $x$ & Multiplicative identity \\
7 & $1 \times x$ & $x$ & Commutative multiplicative identity \\
8 & $x \times 0$ & $0$ & Multiplicative annihilator \\
9 & $0 \times x$ & $0$ & Commutative multiplicative annihilator \\
10 & $x \times (-1)$ & $-x$ & Multiplicative negation \\
11 & $(-1) \times x$ & $-x$ & Commutative multiplicative negation \\
12 & $x / 1$ & $x$ & Divisive identity \\
13 & $x / (-1)$ & $-x$ & Divisive negation \\
14 & $x / x \quad (x \ne 0)$ & $1$ & Multiplicative inverse identity \\
15 & $0 / x \quad (x \ne 0)$ & $0$ & Zero numerator absorption \\
16 & $x \pmod 1$ & $0$ & Unit modulus nullity \\
17 & $x \pmod x \quad (x \ne 0)$ & $0$ & Self-modulus absorption \\
18 & $0 \pmod x \quad (x \ne 0)$ & $0$ & Zero modulus nullity \\
19 & $x \ll 0$ & $x$ & Zero bit-shift nullity \\
20 & $0 \ll k$ & $0$ & Zero shifted absorption \\
21 & $x \gg 0$ & $x$ & Zero arithmetic shift nullity \\
22 & $0 \gg k$ & $0$ & Zero shifted absorption \\
23 & $x \mathbin{\&} x$ & $x$ & Bitwise AND idempotence \\
24 & $x \mathbin{\&} 0$ & $0$ & Bitwise AND annihilator \\
25 & $x \mathbin{\&} (-1)$ & $x$ & Bitwise AND full-mask identity \\
26 & $x \mathbin{|} x$ & $x$ & Bitwise OR idempotence \\
27 & $x \mathbin{|} 0$ & $x$ & Bitwise OR identity \\
28 & $x \mathbin{|} (-1)$ & $-1$ & Bitwise OR full-mask annihilator \\
29 & $x \oplus x$ & $0$ & Bitwise XOR self-inversion \\
30 & $x \oplus 0$ & $x$ & Bitwise XOR identity \\
\bottomrule
\end{tabular}
\caption{The 30 fundamental algebraic identities enforced during term interning.}
\label{tab:algebraic_identities}
\end{table}

\subsection{Formal Correctness Proof for Reassociation}
In addition to identity elimination, the interner enforces canonical constant reassociation:
\begin{align}
(x + c_1) + c_2 &\longrightarrow x + (c_1 + c_2) \\
(x - c_1) + c_2 &\longrightarrow x + (c_2 - c_1) \\
(x \times c_1) \times c_2 &\longrightarrow x \times (c_1 \times c_2)
\end{align}
Because two's-complement arithmetic over fixed-width 64-bit integers forms a commutative ring $(\mathbb{Z} / 2^{64}\mathbb{Z}, +, \times)$, addition and multiplication are strictly associative and commutative. Thus, these rewrites are semantics-preserving under all possible overflow conditions.

\section{Process Trees: Nodes, Edges, and Topology}
\label{sec:symbolic_state:trees}

The execution of the supercompiler constructs an explicit directed hypergraph termed the \emph{Process Tree} ($\proctree$). Nodes in the process tree are categorized into four canonical variants:

\begin{enumerate}
    \item \textbf{Leaf Nodes}: Terminal evaluation states representing either an explicit return value (\texttt{Terminator::Return}), an unreachable state resulting from interval dead-branch pruning, or a runtime divergence.
    \item \textbf{Branch Nodes}: Control-flow decision points where driving bifurcated into multiple child configurations (corresponding to conditional branches or pattern matches).
    \item \textbf{Knot Nodes}: Generalization loop headers where recursive execution has folded into a fixed point configuration. A knot node defines the entry signature for a synthesized recursive residual function.
    \item \textbf{Fold Nodes}: Leaf-level back-edges referencing an ancestor knot node. A fold node supplies the concrete arguments for the recursive invocation.
\end{enumerate}

\section{Knot Allocation and Generalization Fallback}
\label{sec:symbolic_state:knots}

When the whistle detects that a newly reached configuration $C_{\text{curr}}$ homeomorphically embeds an ancestor configuration $C_{\text{anc}}$, the driving engine halts linear evaluation:

\begin{enumerate}
    \item If $C_{\text{curr}}$ is an exact $\alpha$-rename of $C_{\text{anc}}$, a fold node is materialized pointing to $C_{\text{anc}}$, closing the recursive loop.
    \item If $C_{\text{curr}}$ structurally contains $C_{\text{anc}}$ with varying accumulator values, the supercompiler invokes \emph{Most Specific Generalization} (MSG, Chapter~\ref{chap:generalization}).
    \item The MSG algorithm computes the common generalization $C_{\text{gen}} = \text{MSG}(C_{\text{anc}}, C_{\text{curr}})$. The original subtree rooted at $C_{\text{anc}}$ is pruned and replaced by $C_{\text{gen}}$, which is promoted to a Knot node. Driving resumes from $C_{\text{gen}}$.
\end{enumerate}

\section{Residualization: Synthesizing SSA MIR from Trees}
\label{sec:symbolic_state:residualize}

Once all branches of the process tree have terminated at either Leaves or Folds, the supercompiler executes \emph{residualization} (\texttt{src/mir/supercompiler/residualize.rs}):

\begin{enumerate}
    \item \textbf{Function Synthesis}: Each Knot node is allocated a fresh residual function signature $\mathtt{res\_fn\_}k(\vec{p})$.
    \item \textbf{Block Emission}: Tree edges are linearized into SSA basic blocks. Branch nodes emit \texttt{Terminator::BranchIf} or \texttt{Terminator::Switch}.
    \item \textbf{Fold Lowering}: Fold nodes emit direct function calls $\mathtt{Call}(\mathtt{res\_fn\_}k, \vec{args})$, completing the synthesized recurrence.
    \item \textbf{SSA Verification}: The residualized function undergoes standard dominance and SSA validation prior to native backend lowering.
\end{enumerate}

\section{Canonical Hash-Consing and Algebraic Simplification (Phase 59)}
\label{sec:state:algebraic_impl}

Symbolic terms (\texttt{SymTerm}) are interned via \texttt{TermInterner} in \texttt{src/mir/supercompiler/term.rs}. During term interning, the 30 canonical algebraic identities are applied in constant time:

\begin{lstlisting}[language=Rust, caption={Term Representation, Hash-Consing, and Algebraic Reductions (\texttt{src/mir/supercompiler/term.rs}).}, label={lst:term_intern_impl}]
    /// A reference to a heap-allocated value (symbolic address).
    Ref(SymTermId, Type),
    /// A dereference of a symbolic pointer.
    Deref(SymTermId, Type),
    /// The discriminant (variant tag) of an enum value.
    Discriminant(SymTermId, Type),
    /// A symbolically known closure: fn_name + captured symbolic terms.
    ClosureVal(String, Vec<SymTermId>, Type),
    /// An unevaluated thunk: suspended body + captured symbolic terms.
    Thunk(String, Vec<SymTermId>, Type),
}

impl SymTerm {
    pub fn ty(&self) -> &Type {
        match self {
            SymTerm::ConstInt(_, ty) => ty,
            SymTerm::ConstFloat(_, ty) => ty,
            SymTerm::ConstBool(_) => &Type::Bool,
            SymTerm::ConstStr(_) => &Type::Str,
            SymTerm::Var(_, ty) => ty,
            SymTerm::Binary(_, _, _, ty) => ty,
            SymTerm::Unary(_, _, ty) => ty,
            SymTerm::Constructor(_, _, _, ty) => ty,
            SymTerm::Select(_, _, _, ty) => ty,
            SymTerm::Phi(_, ty) => ty,
            SymTerm::Call(_, _, ty) => ty,
            SymTerm::Ref(_, ty) => ty,
            SymTerm::Deref(_, ty) => ty,
            SymTerm::Discriminant(_, ty) => ty,
            SymTerm::ClosureVal(_, _, ty) => ty,
            SymTerm::Thunk(_, _, ty) => ty,
        }
    }

    pub fn to_literal(&self) -> Option<TypedLiteral> {
        match self {
            SymTerm::ConstInt(val, ty) => Some(TypedLiteral::Int(*val, ty.clone())),
            SymTerm::ConstFloat(bits, ty) => {
                Some(TypedLiteral::Float(f64::from_bits(*bits), ty.clone()))
            }
            SymTerm::ConstBool(b) => Some(TypedLiteral::Bool(*b)),
            SymTerm::ConstStr(s) => Some(TypedLiteral::Str(s.clone())),
            _ => None,
        }
    }
}

#[derive(Debug, Clone, Default)]
pub struct TermInterner {
    terms: Vec<SymTerm>,
    lookup: HashMap<SymTerm, SymTermId>,
    sizes: Vec<usize>,
    depths: Vec<usize>,
    hashes: Vec<u64>,
}

const FX_K: u64 = 0x517cc1b727220a95;

#[inline]
fn fx_hash_step(hash: u64, val: u64) -> u64 {
    hash.rotate_left(5) ^ val.wrapping_mul(FX_K)
}

fn fx_hash_bytes(mut hash: u64, bytes: &[u8]) -> u64 {
    for &b in bytes {
        hash = fx_hash_step(hash, b as u64);
    }
    hash
}

fn hash_place(mut h: u64, p: &Place) -> u64 {
    h = fx_hash_bytes(h, p.local.as_bytes());
    for proj in &p.projections {
        match proj {
            crate::mir::Projection::Deref => {
                h = fx_hash_step(h, 0xd37ef);
            }
            crate::mir::Projection::Field(f) => {
                h = fx_hash_bytes(fx_hash_step(h, 0xf1e1d), f.as_bytes());
            }
            crate::mir::Projection::Index(idx) => {
                h = hash_place(fx_hash_step(h, 0x14d3c), idx);
            }
            crate::mir::Projection::Payload(idx) => {
                h = fx_hash_step(fx_hash_step(h, 0x9a710ad), *idx as u64);
            }
        }
    }
    h
}

fn binary_op_discriminant(op: BinaryOp) -> u64 {
    match op {
        BinaryOp::Add => 1,
        BinaryOp::Sub => 2,
        BinaryOp::Mul => 3,
        BinaryOp::Div => 4,
        BinaryOp::Mod => 5,
        BinaryOp::Pow => 6,
        BinaryOp::BitAnd => 7,
        BinaryOp::BitOr => 8,
        BinaryOp::BitXor => 9,
        BinaryOp::Shl => 10,
        BinaryOp::Shr => 11,
        BinaryOp::Eq => 12,
        BinaryOp::Ne => 13,
        BinaryOp::Lt => 14,
        BinaryOp::Le => 15,
        BinaryOp::Gt => 16,
        BinaryOp::Ge => 17,
    }
}

fn unary_op_discriminant(op: UnaryOp) -> u64 {
    match op {
        UnaryOp::Neg => 1,
        UnaryOp::Not => 2,
    }
}

impl TermInterner {
    pub fn new() -> Self {
        TermInterner {
            terms: Vec::new(),
            lookup: HashMap::new(),
            sizes: Vec::new(),
            depths: Vec::new(),
            hashes: Vec::new(),
        }
    }

    pub fn get(&self, id: SymTermId) -> &SymTerm {
        &self.terms[id.0]
    }

    pub fn len(&self) -> usize {
        self.terms.len()
    }

    pub fn is_empty(&self) -> bool {
        self.terms.is_empty()
    }

    pub fn size(&self, id: SymTermId) -> usize {
        self.sizes[id.0]
    }

    pub fn depth(&self, id: SymTermId) -> usize {
        self.depths[id.0]
    }

    pub fn hash(&self, id: SymTermId) -> u64 {
        self.hashes[id.0]
    }

    #[inline]
    pub fn structural_eq(&self, t1: SymTermId, t2: SymTermId) -> bool {
        t1 == t2
    }

    pub fn intern_const(&mut self, lit: TypedLiteral) -> SymTermId {
        match lit {
            TypedLiteral::Int(i, ty) => self.intern(SymTerm::ConstInt(i, ty)),
            TypedLiteral::Float(f, ty) => self.intern(SymTerm::ConstFloat(f.to_bits(), ty)),
            TypedLiteral::Bool(b) => self.intern(SymTerm::ConstBool(b)),
            TypedLiteral::Str(s) => self.intern(SymTerm::ConstStr(s)),
        }
    }

    pub fn intern_int(&mut self, val: i64) -> SymTermId {
        self.intern(SymTerm::ConstInt(val, Type::I64))
    }

    pub fn intern_bool(&mut self, val: bool) -> SymTermId {
        self.intern(SymTerm::ConstBool(val))
    }

    pub fn intern_float(&mut self, val: f64) -> SymTermId {
        self.intern(SymTerm::ConstFloat(val.to_bits(), Type::F64))
    }

    pub fn intern_typed_zero(&mut self, ty: &Type) -> SymTermId {
        match ty {
            Type::F64 | Type::F32 => self.intern(SymTerm::ConstFloat(0.0f64.to_bits(), ty.clone())),
            Type::Bool => self.intern(SymTerm::ConstBool(false)),
            _ => self.intern(SymTerm::ConstInt(0, ty.clone())),
        }
    }

    pub fn intern_typed_one(&mut self, ty: &Type) -> SymTermId {
        match ty {
            Type::F64 | Type::F32 => self.intern(SymTerm::ConstFloat(1.0f64.to_bits(), ty.clone())),
            Type::Bool => self.intern(SymTerm::ConstBool(true)),
            _ => self.intern(SymTerm::ConstInt(1, ty.clone())),
        }
    }

    pub fn intern_var(&mut self, place: Place, ty: Type) -> SymTermId {
        self.intern(SymTerm::Var(place, ty))
    }

    pub fn intern_binary(
        &mut self,
        op: BinaryOp,
        mut left: SymTermId,
        mut right: SymTermId,
        ty: Type,
    ) -> SymTermId {
        let mut left_term = self.get(left).clone();
        let mut right_term = self.get(right).clone();

        // 1. Constant folding
        if let (Some(l_lit), Some(r_lit)) = (left_term.to_literal(), right_term.to_literal()) {
            if let Some(folded) = fold_const_binary(op, &l_lit, &r_lit) {
                return self.intern_const(folded);
            }
        }

        // 2. Canonical commutative reordering (ALG-01)
        // Constants are moved to the right; non-constants are sorted deterministically by SymTermId.
        if is_commutative(op) {
            let should_swap = match (is_const(&left_term), is_const(&right_term)) {
                (true, false) => true,
                (false, true) => false,
                _ => left.0 > right.0,
            };
            if should_swap {
                std::mem::swap(&mut left, &mut right);
                std::mem::swap(&mut left_term, &mut right_term);
            }
        }

        // 3. Structural algebraic simplifications (ALG-02, ALG-04)
        match op {
            BinaryOp::Add => {
                if is_zero(&right_term) {
                    return left;
                }
                if is_zero(&left_term) {
                    return right;
                }
                // Constant reassociation: (x + c1) + c2 => x + (c1 + c2)
                if let SymTerm::Binary(BinaryOp::Add, inner_x, inner_c, _) = left_term {
                    if let (Some(c1_lit), Some(c2_lit)) = (self.get(inner_c).to_literal(), right_term.to_literal()) {
                        if let Some(c_sum) = fold_const_binary(BinaryOp::Add, &c1_lit, &c2_lit) {
                            let c_sum_id = self.intern_const(c_sum);
                            return self.intern_binary(BinaryOp::Add, inner_x, c_sum_id, ty);
                        }
                    }
                }
                // Constant reassociation: (x - c1) + c2 => x + (c2 - c1)
                if let SymTerm::Binary(BinaryOp::Sub, inner_x, inner_c, _) = left_term {
                    if let (Some(c1_lit), Some(c2_lit)) = (self.get(inner_c).to_literal(), right_term.to_literal()) {
                        if let Some(c_diff) = fold_const_binary(BinaryOp::Sub, &c2_lit, &c1_lit) {
                            let c_diff_id = self.intern_const(c_diff);
                            return self.intern_binary(BinaryOp::Add, inner_x, c_diff_id, ty);
                        }
                    }
                }
            }

\end{lstlisting}
